\section{Evaluation}

\label{sec:evaluation}

\subsection{Deployment}

We deployed ZooCoord across 14 research groups in the Zoo Network, spanning three scientific domains:

\begin{table}[H]
\centering
\caption{Deployment overview by research domain.}
\label{tab:deployment}
\begin{tabularx}{\textwidth}{@{}lcccX@{}}
\toprule
\textbf{Domain} & \textbf{Groups} & \textbf{Researchers} & \textbf{Experiments} & \textbf{Description} \\
\midrule
Biodiversity Monitoring & 6 & 47 & 412 & Camera trap analysis, species surveys, population estimates \\
Ecological Modeling & 5 & 38 & 287 & Species distribution, climate impact, habitat connectivity \\
Conservation Genetics & 3 & 22 & 148 & eDNA analysis, population genetics, phylogeography \\
\midrule
\textbf{Total} & 14 & 107 & 847 & \\
\bottomrule
\end{tabularx}
\end{table}

\subsection{Reproducibility Outcomes}

Of 847 committed experiments, 773 reached completion (91.3\%).
The remaining 74 were withdrawn (42), suspended pending additional data (27), or abandoned (5).

\begin{table}[H]
\centering
\caption{Reproducibility verification results.}
\label{tab:reproducibility}
\begin{tabular}{@{}lcccc@{}}
\toprule
\textbf{Domain} & \textbf{Verified} & \textbf{Reproduced} & \textbf{Partial} & \textbf{Failed} \\
\midrule
Biodiversity Monitoring & 289 & 271 (93.8\%) & 14 (4.8\%) & 4 (1.4\%) \\
Ecological Modeling & 198 & 176 (88.9\%) & 16 (8.1\%) & 6 (3.0\%) \\
Conservation Genetics & 104 & 89 (85.6\%) & 12 (11.5\%) & 3 (2.9\%) \\
\midrule
\textbf{Total} & 591 & 536 (90.7\%) & 42 (7.1\%) & 13 (2.2\%) \\
\bottomrule
\end{tabular}
\end{table}

The 13 reproduction failures were investigated through dispute resolution.
Seven resulted from undocumented environment dependencies (now captured by enhanced provenance logging), four from genuine methodological errors caught before publication, and two from hardware-specific numerical instabilities.

\subsection{Time-to-Reproduction}

We measured the time required for an independent party to reproduce an experiment's results:

\begin{table}[H]
\centering
\caption{Time-to-reproduction comparison (median days).}
\label{tab:time}
\begin{tabular}{@{}lcc@{}}
\toprule
\textbf{Complexity} & \textbf{Traditional} & \textbf{ZooCoord} \\
\midrule
Simple (1--5 FEG nodes) & $14.2 \pm 8.7$ & $0.3 \pm 0.2$ \\
Medium (6--20 nodes) & $42.6 \pm 21.3$ & $2.1 \pm 1.4$ \\
Complex (21+ nodes) & $127.4 \pm 64.8$ & $8.7 \pm 5.2$ \\
\midrule
\textbf{Overall} & $48.3 \pm 41.2$ & $3.2 \pm 4.1$ \\
\bottomrule
\end{tabular}
\end{table}

ZooCoord reduced median time-to-reproduction by 93.4\% overall.
The improvement is most dramatic for simple experiments, where fully automated re-execution requires only hours.
Complex experiments still require human intervention for data access authorization and environment-specific configuration.

\subsection{Data Sharing}

Data sharing rates increased substantially under ZooCoord's incentivized framework:

\begin{table}[H]
\centering
\caption{Data sharing metrics before and after ZooCoord deployment.}
\label{tab:sharing}
\begin{tabular}{@{}lcc@{}}
\toprule
\textbf{Metric} & \textbf{Pre-ZooCoord} & \textbf{Post-ZooCoord} \\
\midrule
Experiments with shared raw data & 18.3\% & 74.2\% \\
Experiments with shared code & 42.7\% & 97.1\% \\
Experiments with shared intermediate results & 7.2\% & 81.6\% \\
Cross-group data reuse events & 4 per month & 17 per month \\
\bottomrule
\end{tabular}
\end{table}

\subsection{Deviation Detection}

The commitment protocol detected 23 potential reproducibility issues during the execution phase, before results were finalized:

\begin{table}[H]
\centering
\caption{Deviation types detected during experiment execution.}
\label{tab:deviations}
\begin{tabular}{@{}lcc@{}}
\toprule
\textbf{Deviation Type} & \textbf{Count} & \textbf{Severity} \\
\midrule
Undeclared parameter changes & 8 & Material \\
Sample size deviations & 6 & Registered \\
Analysis script modifications & 5 & Material \\
Unreported data exclusions & 3 & Commitment Violation \\
Software version drift & 1 & Registered \\
\bottomrule
\end{tabular}
\end{table}

All three commitment violations were resolved through revised commitments with transparent documentation of the reasons for deviation.
This early detection mechanism prevented potential reproducibility failures from propagating to published results.

\subsection{Cross-Institutional Collaboration}

ZooCoord facilitated 12 cross-institutional collaborations that emerged organically from the shared experiment registry:

\begin{itemize}[leftmargin=*]
    \item 5 involved combining datasets from multiple sites for meta-analyses.
    \item 4 involved one group extending another group's methodology to a new geographic region.
    \item 3 involved verifiers discovering methodological improvements during reproduction and co-authoring enhanced versions.
\end{itemize}

These collaborations produced 8 joint publications and 3 shared datasets, none of which would have occurred under traditional coordination mechanisms where experiment details are typically shared only at publication.

% ============================================================
% 9. Economic Analysis
% ============================================================
